% part: intuitionistic-logic
% chapter: semantics
% section: introduction

\documentclass[../../../include/open-logic-section]{subfiles}

\begin{document}

\olfileid{int}{sem}{int}

\olsection{引言}

没有语义学，任何逻辑都无法得到令人满意的描述，直觉主义逻辑也不例外。对于经典逻辑，基于赋值的语义学是典范的；而直觉主义逻辑则有几种相互竞争的语义学。但它们都算不上彻底令人满意，因为它们都没能以直觉主义可接受的方式给出联结词含义的说明。

基于关系模型的语义学或许是最流行的一种，它类似于模态逻辑的语义学。在这种语义学中，命题变元被指派给各个世界，而这些世界之间由一个可达关系相联系。该关系总是一个偏序，即它是自反的、反对称的和传递的。

直观上，你可能会把这些世界理解为知识状态或“证据情境”。当且仅当就我们所知，$w'$ 是一个可能的（未来的）知识状态时，状态 $w'$ 从 $w$ 才是可达的，也就是说，$w'$ 与在 $w$ 处所知的内容相容。一旦一个命题被知道，它就不可能再变为未知，即如果在 $w$ 处知道 $!A$ 且 $Rww'$，那么在 $w'$ 处也知道 $!A$。因此，“知识”相对于可达关系是单调的。

如果我们像认知逻辑那样，将“知道 $!A$”定义为“在所有认知替代世界中都为真”，那么 $!A \land !B$ 在 $w$ 处被知道，当且仅当在所有认知替代世界中，$!A$ 和 $!B$ 都被知道。但由于知识是单调的且 $R$ 是自反的，这就意味着 $!A \land !B$ 在 $w$ 处被知道，当且仅当 $!A$ 和 $!B$ 在 $w$ 处被知道。出于同样的原因，$!A \lor !B$ 在 $w$ 处被知道，当且仅当至少有一个在 $w$ 处被知道。因此，对于 $\land$ 和 $\lor$，联结词的真值条件与经典逻辑中的相同。

然而，条件句的真值条件不同于经典逻辑。$!A \lif !B$ 在 $w$ 处被知道，当且仅当不存在满足 $Rww'$ 的世界 $w'$，使得在 $w'$ 处 $!A$ 被知道而 $!B$ 不被知道。这不同于“$!A$ 未知或 $!B$ 在 $w$ 处被知道”这一条件。因为如果我们在 $w$ 处既不知道 $!A$ 也不知道 $!B$，那么可能存在一个未来的认知状态 $w'$（$Rww'$），使得在 $w'$ 处，我们知道了 $!A$ 而没有随之知道 $!B$。

我们知道 $\lnot !A$，仅当不存在任何可能的未来认知状态使我们在其中知道 $!A$。这里的想法是，如果 $!A$ 是可知的，那么在某个可能的未来认知状态中，$!A$ 就会变为已知的。由于我们不可能知道 $\lfalse$，在那个未来认知状态中，我们将知道 $!A$ 却不知道 $\lfalse$。

在这种解释下，排中律失效。因为存在一些我们尚不知道、但未来可能知道的 $!A$。对于这样的公式 $!A$，$!A$ 和 $\lnot !A$ 都是未知的，所以 $!A \lor \lnot !A$ 不被知道。但是，我们的确知道，例如 $\lnot(!A \land \lnot !A)$。因为不存在任何可能的未来状态能让我们同时知道 $!A$ 和 $\lnot !A$，而且这一点与我们是否知道 $!A$ 或 $\lnot !A$ 无关。

关系模型并不是直觉主义逻辑唯一可用的语义学。拓扑语义学是另一种：在这里，命题被解释为拓扑空间中的开集，而联结词被解释为对这些集合的运算（例如，$\land$ 对应于交集）。

\end{document}